% ===== CHAPTER 4: CONCLUSION =====

\chapter{Conclusion}

\section{Summary of Work}

TravelSync set out to replace the fragmented, multi-application process of
planning a trip with a single platform. The team built a React single-page
application backed by an asynchronous FastAPI service and MongoDB, organised
around the idea that a chosen destination becomes a trip and every
supporting feature works from that trip.

All four objectives from Section~\ref{sec:objectives} were achieved. An undecided traveller
receives ranked, explained destination recommendations from a weighted
nearest-neighbour engine over 250 destinations that respects duration,
budget and distance constraints (Objective~1). Each trip opens a workspace
with its top places, a two-tier weather-aware packing list, scored
accommodation recommendations, historical health context and a map with
driving distances (Objective~2). Groups collaborate through invite codes,
live WebSocket chat, structured suggestions, one-member-one-vote voting and
an owner-confirmed decision (Objective~3). Bookings with private file
attachments, trip history, notifications and a privacy-preserving assistant
complete the platform, protected throughout by session security and per-user
authorisation (Objective~4). The system is covered by 303 passing automated
tests.

\section{Key Findings}

\begin{itemize}
  \item \textbf{Finding 1:} A transparent, content-based nearest-neighbour
    ranking combined with simple constraint scores was sufficient for
    destination recommendation without any training data. The engine
    reached a Precision@5 of 97.5\% against a random baseline of 57\%,
    answered each request in under 2\,ms, and never violated the budget
    ceiling or the distance tolerance.
  \item \textbf{Finding 2:} Designing every external dependency to fail
    gracefully, through caching, per-user limits and local fallbacks, made
    it possible to build on free-tier services while keeping the application
    usable when a provider was slow, exhausted or unavailable.
  \item \textbf{Finding 3:} The evaluation exposed an unexpected interaction
    in the ranking: the fixed weight given to the traveller type can
    outweigh a single selected interest, so that for some combinations (for
    example Culture with the Solo type) the recommendations favour the
    traveller-type profile over the interest itself.
\end{itemize}

\section{Future Work and Recommendations}
\label{sec:future-work}

The current version provides a complete foundation that can be extended in
the following directions:

\begin{enumerate}
  \item Potential enhancements to the system: tune the traveller-type
    weight and the final-score weights using recorded user choices, so that
    the recommender learns from which destinations users actually select
    \cite{lops2011content}.
  \item Alternative approaches worth exploring: combine the content-based
    engine with collaborative filtering once enough trips have been
    recorded, and replace the hand-curated health alerts with an automated
    feed from official surveillance sources \cite{idsp}.
  \item Scalability improvements: move the in-process chat hub to a shared
    message broker so that several backend workers can serve WebSocket
    clients, and move booking files to object storage as their volume grows.
  \item Integration with other systems: link TravelSync with transport and
    hotel providers that publish real prices and availability, and allow
    itineraries to be exported to calendar applications.
  \item Real-world deployment considerations: conduct a formal usability
    study, add multilingual support, and complete a production deployment
    with HTTPS, monitoring and automated backups of the MongoDB database.
\end{enumerate}

\section{Final Remarks}

Building TravelSync gave the team practical experience of the complete
software engineering life cycle: turning a proposal into requirements and
UML models, revising the design when implementation revealed better options
(such as replacing a planned administrator role with versioned datasets, and
a separate messaging service with native WebSockets), and validating the
result with automated tests and measurements. The most valuable lessons
were the importance of explainable algorithms that can be defended to users
and evaluators, of security and authorisation designed in from the first
increment rather than added at the end, and of honest presentation of data,
such as labelling estimated prices and historical health information for
what they are. The resulting platform demonstrates that discovery,
preparation and group coordination for a trip can be brought together in one
coherent and trustworthy experience.
